\documentclass[border=8pt]{standalone}
\usepackage[american]{circuitikz}
\begin{document}
\begin{circuitikz}[font=\sffamily]
\node[op amp] (A) at (5,2.2) {};
\draw (A.out)--++(1,0) coordinate(vo) node[right]{$v_o=\pm V_S$};
\draw (A.+)--++(-.25,0) coordinate(vp);
\draw (vp)|-(3.1,4.5) to[R,l=$R_1$] (5.7,4.5)-|(vo);
\draw (vp) to[R,l_=$R_2$] ++(0,-2.0) node[ground]{};
\draw (A.-)--++(-1.2,0) coordinate(trig);
\draw (trig) to[R] ++(0,-2.25) node[ground]{};
\draw (trig) to[C] ++(-1.5,1.2) node[left,align=right]{SET\\negative pulse};
\draw (trig) to[C] ++(-1.5,-1.2) node[left,align=right]{RESET\\positive pulse};
\node at (1.72,3.28) {$C_S$};
\node at (1.22,2.10) {$C_R$};
\node[right] at (2.93,1.43) {$R_B$};
\node[above] at (4.4,5.3) {Regenerative bistable memory};
\node[align=center] at (4.6,-1.25) {state persists after the trigger is removed};
\end{circuitikz}
\end{document}
